\section{Conclusão}

As séries de Fourier foram verificadas nos 24 cenários, com sete
pares de coeficientes por caso. A FFT das amostras geradas no
Simulink apresentou erro máximo de $\ErroCoefMax$ V por par;
os erros relativos máximos da fundamental e da energia total
foram $\ErroCUmMax\%$ e $\ErroEnergiaMax\%$, respectivamente.
A concordância foi conferida também por integração direta das
ondas e pelo teorema de Parseval.

O triângulo bipolar concentra $98{,}5534\%$ de sua energia na
fundamental e $99{,}9689\%$ até a ordem 7. Na rampa unipolar,
a reconstrução com a média e as ordens 1 a 7 retém $97{,}9766\%$;
a média sozinha responde por $75\%$ do total. O
$\mathrm{THD}_7$ de $11{,}9842\%$ dos triângulos, frente a
$71{,}5400\%$ das rampas, quantifica a menor participação relativa
dos harmônicos superiores nas ondas triangulares.

Por fim, notou-se que ao deslocar a onda a amplitude é preservada e as razões de energia também.
Reduzir o período para $0{,}8P$ eleva a frequência de 60 para 75 Hz;
alternar dois ciclos distintos produz a fundamental de 30 Hz e
participações de $6{,}6934\%$ e $11{,}1557\%$ na energia dos
retangulares unipolar e bipolar. Os pequenos erros da FFT são
compatíveis com a discretização adotada; as oscilações junto aos
saltos nas reconstruções decorrem do truncamento da série.
